

\section{The agent's own reading}\label{sec:reading}

\emph{This section is the agent's opinion. The sections before it are results. You may disagree with this one.}

\begin{itemize}
\item At 8 bits the claim failed, but the test was not strict. Every new creature added a row of
  \N{rule.columns} trained numbers, so the model grew with the vocabulary. I think that is the most
  likely reason the facts were all known, but it is a guess. The frozen-table test is the way to find out.
\item The claim began to hold up where the bits per creature were high. At 64 bits and with many
  creatures the scores fell, and a bigger body helped only a little. This looks more like a limit of
  room than of time, but the curves were still rising, so it may be time. I cannot tell.
\item ``not'' known fell much further than facts known, down to \N{bits64.notk.lowest} at 64 bits.
  A statement about what a creature does \emph{not} do may be harder to learn than one about what it
  does. This was not tested.
\item Unequal use made everything worse at 64 bits. The creatures at the end of the list are met
  rarely, so they may only need more training.
\item My expectation in P4 was half right: nothing fell with more creatures at 8, 16 and 32 bits,
  and something did at 64 bits.
\end{itemize}

\section{Limits, and what this does not show}\label{sec:limits}

\textbf{Limits.}
\begin{itemize}
\item \textbf{The model grew with the vocabulary.} The word table was trained, so every new word
  added a row of \N{rule.columns} trained numbers (figure \ref{fig:params}). The claim, read
  strictly, needs a model of fixed size. This test is the case you chose, not the strict one.
\item \textbf{Few repeats where it matters most.} 8 bits has \N{rule.repeats} repeats, 16 and 32 bits
  have 3, 64 bits has 2. Two repeats show a spread but cannot say how often a result happens.
\item \textbf{One kind of fact.} Every fact is a link from a creature to one value. Facts that
  link creatures to each other were not tested.
\item \textbf{A fixed length.} Every setting of a stage trained for the same number of steps. A
  setting that was slower for a good reason is judged at the same time as the others.
\item \textbf{Chance in the exam.} Questions about new combinations with 16 creatures rest on
  only \N{qnew.count.c16} questions. The score is coarse there.
\item \textbf{Two places.} Part of the 16 and 32-bit stages ran on your computer and the rest on
  Modal. The same programs ran in both places. I did not check whether a run gives the same digits in
  both.
\end{itemize}

\textbf{What this does not show.}
\begin{itemize}
\item It does not show \emph{why} the score stayed at 100 at 8 bits. The extra storage of each new
  word is one explanation. It was not tested.
\item It does not show that a model of fixed size can hold any number of words. The frozen-table
  test, where a new word brings no storage of its own, is ready and \textbf{has not been run}.
\item It does not tell ``cannot hold'' from ``not learned yet'' for any setting whose score was
  still rising when training stopped, because the plan's longer control run (four times as long, for
  each setting that fails) was not done.
\item It does not say what happens with more than 64 bits, with other model designs, or in real
  language, where the meaning of a word depends on other words.
\end{itemize}

\section{Everything done differently from the plan}\label{sec:extras}

\begin{enumerate}[itemsep=1.5pt]
\item \textbf{The 64-bit stage was cut} from the plan's full grid with three repeats (120 runs) to
  \N{expected.bits64} runs: 16, 256 and 1,024 creatures, body widths 8, 32 and 128, both kinds of
  use, 2 repeats. This was to stay inside your \$30 of cloud credit.
\item \textbf{Part of the 16 and 32-bit stages ran on your computer,} for the same reason:
  \N{local.bits16} runs of the 16-bit stage and \N{local.bits32} of the 32-bit stage (repeat 2, for
  some body widths), and the practice runs of stage \texttt{pilot2}
  (\N{expected.pilot2} runs). The rest ran on Modal.
\item \textbf{A value word was replaced.} \s{mice} was also a creature name, so the value
  \s{mice} was replaced by \s{gnats}.
\item \textbf{Body widths up to 64 trained on 1 thread and width 128 on 2.} This changes the cost,
  not the method.
\item \textbf{The longer control run was not done.} The plan said one run four times as long for
  each setting that fails. None was run.
\item \textbf{Extra runs.} A small timing stage of 4 runs (\texttt{speed}) measured the cloud
  machines. A tiny end-to-end test of 1 run (\texttt{smoke}) checked the programs.
\item \textbf{Settings per stage.} The learning rate and the length were fixed from practice runs
  separately for 8 bits, for 64 bits, and for 16 and 32 bits (section \ref{sec:practice}), each
  fixed before the counted runs of that stage started.
\item \textbf{Missing runs.} None. All \N{runs.main8}, \N{runs.bits16}, \N{runs.bits32} and
  \N{runs.bits64} counted runs of the four stages were done and read (planned: \N{expected.main8},
  \N{expected.bits16}, \N{expected.bits32} and \N{expected.bits64}). Some runs were interrupted and went on from their saved checkpoint, which the programs are built to do exactly (8 bits: \N{resumed.main8}; 16 bits: \N{resumed.bits16}; 32 bits: \N{resumed.bits32}; 64 bits: \N{resumed.bits64}). Each is counted once.
\end{enumerate}

\textbf{Where it ran, and what it cost.} Modal is a cloud service. Up to 100 small machines ran at
once, with 1 or 2 processor cores and no graphics card. Together the counted runs took about
\N{hours.main8} machine hours (8 bits), \N{hours.bits16} (16 bits), \N{hours.bits32} (32 bits) and
\N{hours.bits64} (64 bits), counting runs on all machines. Modal's own cost meter read
\$\N{cost.meter} at the end (\N{cost.meter.source}); it read \$\N{cost.log.first} when
\N{cost.log.first.time} was first logged and \$\N{cost.log.last} at \N{cost.log.last.time}, the last
logged reading.

\section{What could come next}\label{sec:next}

\begin{itemize}
\item \textbf{The frozen-table test.} The same grid as 8 bits (stage \texttt{frozen8}, 200 runs),
  with a table that is never trained, so that a new word brings no storage of its own. Before it
  was run, the agent's expectation was: body width 8 holds 64 creatures but not 256, width 32
  holds 256 but not 1,024, and width 128 holds 1,024. It is ready and has not been run.
\end{itemize}

The open questions in \texttt{research\_questions.md}, in a line each:
\begin{itemize}
\item \textbf{Q1:} Can a model say ``no'' about creatures it was never asked a ``no'' question about?
\item \textbf{Q2:} What does capacity depend on: word count, total bits, or bits per number?
  (Sections \ref{sec:bits} and \ref{sec:main8} give a first look.)
\item \textbf{Q3:} Store or query: can a fact learned as a statement be used in a question?
\item \textbf{Q4:} Does the vocabulary matter when words depend on each other?
\item \textbf{Q5:} What changes when a word carries more than 8 bits? (Answered here for 16, 32 and 64
  bits with the plain kind of sentence.)
\item \textbf{Q6:} Does a dimension of its own for every word remove the effect of vocabulary size?
\item \textbf{Q7:} Can rewards teach a model to answer questions where ordinary training did not?
\item \textbf{Q8:} Does more variety in the text give better learning than repeating the same lines?
\item \textbf{Q9:} How can training get out of a stall sooner?
\item \textbf{Q10:} What is the smallest body width that works at all, and was width 8 on 256
  animals at 32 bits too small or only slow?
\item \textbf{Q11:} How does the model hold the facts inside? (Your own study.)
\item \textbf{Q12:} Does it matter whether a line of training text holds one sentence or many?
\end{itemize}
